\section{Updated Conclusions from the Expanded Evaluation}
\label{sec:expanded-final-conclusions}

This completed expansion changes the interpretation from a small-cohort
separation exercise to a source-transfer test with explicit failure cases.
All selected recordings are retained on disk, all four feature families have
item-level completion or observability accounting, and the complete fifteen-by-five
ablation is reproducible from frozen outputs. Spectral evidence was not restricted
to the initial 1,500 AI recordings: the preceding development table already used
7,100 inputs, and the current broad union contains 10,141 recording identities.

The most defensible candidates are S+D for ten-second inputs and S+D+R for the
duration-qualified thirty-second cohort. Their mean source-transfer AUCs are
0.6334 and 0.7641, respectively, with fixed-threshold balanced accuracies of
0.5971 and 0.6986. These are development selection summaries, not the accuracy
of a universal detector. The duration comparison is confounded by a substantial
change in sources and recordings.

The high-frequency hypothesis has measurable value, particularly when the Human
reference distribution is expanded. Nevertheless, greater diversity does not
uniformly improve every transfer direction: the ten-second leader gains on
generator holdout but loses on Human-source holdout. Higher unrestricted-band
scores remain compatible with encoding and separator artifacts. D adds useful
classification information without establishing MIDI velocity recovery. R adds
information in the long cohort but its earlier external rhythm-change validation
still failed. P is unobservable in the short cohort and nearly chance alone in
the long cohort; these measurements cannot support a phrase-creativity claim.

The most consequential negative result is the remaining false-positive burden.
Thirty-second Human specificity is 0.9515 on MusicNet but 0.5233 on DEAM, while
unseen-source Udio classification remains near chance. The pipeline should
therefore be described as an exploratory, interpretable source-sensitive score,
not a conclusive authorship test. A single high score must not be presented as
proof that a recording or artist used generative AI.

Further work should validate rhythm and phrase measurements on longer,
independently annotated music before interpreting those constructs; collect
stronger creator/composition identities for grouping; and reserve new,
genre- and production-diverse Human recordings plus provenance-verified AI
sources for the next frozen evaluation. Any threshold or calibration revision
must use separate validation data, not the locked catalogues just examined.
These are prospective requirements, not additional experiments claimed as
completed in this thesis archive.
